% 1977 年广东卷第 5 题：小山上的电视发射塔 AB（高 50 米），自山下地面 C 点测得
% 塔底 B 仰角 40°、塔顶 A 仰角 70°，D 在 B 正下方，右侧标 50 米尺寸线（按源卷第 3 页布局）
\begin{tikzpicture}[scale=0.75]
  % 地面
  \draw[thick] (-0.8,0) -- (9.8,0);
  % 小山轮廓
  \draw[thick] (1.6,0) .. controls (3.0,1.3) and (4.6,2.3) .. (6,3)
    .. controls (7.2,2.6) and (8.4,1.2) .. (9.4,0);
  \draw[thin] (3.6,0.9) .. controls (4.2,1.5) and (4.8,1.9) .. (5.4,2.4);
  \draw[thin] (6.6,2.4) .. controls (7.2,1.9) and (7.8,1.2) .. (8.2,0.5);
  % 发射塔（梯形塔身加交叉撑）
  \draw[thick] (5.75,3) -- (5.9,5.3) -- (6.1,5.3) -- (6.25,3) -- cycle;
  \draw[thin] (5.79,3.3) -- (6.21,3.9);
  \draw[thin] (6.21,3.3) -- (5.79,3.9);
  \draw[thin] (5.80,3.9) -- (6.20,4.5);
  \draw[thin] (6.20,3.9) -- (5.80,4.5);
  \draw[thin] (5.82,4.5) -- (6.18,5.1);
  \draw[thin] (6.18,4.5) -- (5.82,5.1);
  % 视线（虚线）
  \draw[dashed] (0,0) -- (6,5.3);
  \draw[dashed] (0,0) -- (6,3);
  \draw[dashed] (6,3) -- (6,0);
  % C 处两个仰角弧
  \draw (1.3,0) arc (0:26.57:1.3);
  \draw (1.85,0) arc (0:41.4:1.85);
  \node at (1.5,0.32) {$40^\circ$};
  \node at (1.62,1.05) {$70^\circ$};
  % 50 米尺寸线
  \draw[thin] (6.15,5.3) -- (8.7,5.3);
  \draw[thin] (6.35,3) -- (8.7,3);
  \draw[<->] (8.35,3) -- (8.35,5.3);
  \node[rotate=90] at (8.05,4.15) {50米};
  % 顶点标注
  \fill (6,5.3) circle (1.2pt) node[above] {$A$};
  \fill (6,3) circle (1.2pt) node[right=2pt] {$B$};
  \fill (0,0) circle (1.2pt) node[below] {$C$};
  \fill (6,0) circle (1.2pt) node[below] {$D$};
\end{tikzpicture}
